\section{任务目标、基本原理与可行性}
\label{cr:objectives}
\subsection{课程任务与研究对象}
课程要求理解DRAM、单粒子效应和AI基本方法，在个人计算机与远程计算资源之间
搭建仿真环境，完成基础模型、自然语言改参、Geant4到TCAD传递、多参数扫描以及
拟合预测验证，并由2至3名学生合作完成设计、演示和总结\cite{course}。
据此，项目将目标分解为三个可检查的问题：能否把一个明确的粒子条件转换为
可审查的器件源项；能否在相同电学前提下比较不同尺寸和粒子条件的响应；能否
用部分已求解工况预测尚未参与拟合的工况，再由真实TCAD结果检验预测。

本报告选择1T1C DRAM作为对象。一个访问晶体管控制存储节点与位线的连接，
存储电容上的电荷表征信息。保持时访问管关断，读出时字线打开访问管，
存储电容与预充位线发生电荷共享，感测电路据微小差分判别位值\cite{keeth2001}。
因此，“粒子造成多少电荷变化”和“读出是否仍有足够裕量”是相关而不同的问题。
项目既提取保持段电荷扰动，也建立读出指标，避免仅凭存储节点电压下降量判断
完整存储单元的逻辑状态。

\subsection{单粒子效应与数值描述}
载能粒子穿过半导体时产生电子空穴对，载流子在内部电场中漂移，并发生扩散、
复合及端口收集。当扰动到达敏感节点，可能形成瞬态电流或改变储存电荷，
进一步影响电路判读\cite{label1996}。单粒子瞬态、判据上的错误读出与锁存后
持续存在的位翻转并不等价。本项目尚无真实感放锁存和恢复电路，故主要研究
电荷扰动及明确采样时刻的判据裕量。

在漂移扩散框架中，粒子作用以产生率$G$进入电子、空穴连续性方程：
\begin{equation}
\frac{\partial n}{\partial t}=\frac{1}{q}\nabla\!\cdot\mathbf{J}_n+G-R,\qquad
\frac{\partial p}{\partial t}=-\frac{1}{q}\nabla\!\cdot\mathbf{J}_p+G-R.
\end{equation}
这些方程与泊松方程及外部电路方程耦合，决定存储节点和位线的瞬态变化。
Geant4负责粒子输运及能量沉积，TCAD负责沉积之后的半导体电学响应。
二者之间的接口必须明确材料、坐标、单位、源项形状及有效厚度；工具同时运行
并不足以证明它们描述了相同的物理对象。

\subsection{可行性与实际设计过程}
实际环境由Windows本机承担请求解析、批次管理、归档和后处理，由已有Linux
虚拟机运行Geant4与Sentaurus。本机通过严格SSH入口访问虚拟机，保留命令、
返回状态和输入输出摘要。这实现了远程访问式工作流，但现有材料没有提供
实验室独立服务器或正式资源交接记录，不能把虚拟机使用记录替代课程要求中的
资源交接。现有软件、许可证和脚本已能支持单核串行求解，主要成本来自TCAD
瞬态和必要复验，而非表格或绘图。

设计采用逐级建立证据的方式。v1先用小矩阵贯通真实输运、器件响应与外推；
v2加入统一初态、0/1配对和读出判据；D1进一步追查理想参考留下的读出偏移；
随后R50-main以匹配中性参考及统一预充版本建立更完整的矩阵。
这种过程让前一阶段暴露的问题成为后一阶段的设计依据。
当前报告以共享证据快照为截止点，完成状态见表\ref{cr:progress}；
历史阶段保留其原条件，不用新模型重写旧结论\cite{project}。
